% Source: https://artofproblemsolving.com/community/c3h1693937_divisors
% Stable identity: aops-topic-1693937-q01
Count the pairs $(k,n)$ with $1\le n\le100$ and $k\mid n$ in two ways. For fixed $n$ there are $d(n)$ such $k$; for fixed $k$ the multiples of $k$ up to $100$ number $\left\lfloor\dfrac{100}{k}\right\rfloor$. Hence
\[
\sum_{n=1}^{100}d(n)=\sum_{k=1}^{100}\left\lfloor\dfrac{100}{k}\right\rfloor .
\]
Group the values of $k$:
\begin{itemize}
\item $k=1,\dots,10$: $100+50+33+25+20+16+14+12+11+10=291$;
\item $k=11,\dots,100$: $\left\lfloor\dfrac{100}{k}\right\rfloor=q$ for the $k$ with $\dfrac{100}{q+1}<k\le\dfrac{100}{q}$, i.e.\ $q=9$: $k=11$ ($1$ value); $q=8$: $k=12$ ($1$); $q=7$: $k=13,14$ ($2$); $q=6$: $k=15,16$ ($2$); $q=5$: $k=17,\dots,20$ ($4$); $q=4$: $k=21,\dots,25$ ($5$); $q=3$: $k=26,\dots,33$ ($8$); $q=2$: $k=34,\dots,50$ ($17$); $q=1$: $k=51,\dots,100$ ($50$). This gives $9+8+14+12+20+20+24+34+50=191$.
\end{itemize}
So $\sum d(n)=291+191=482$, and the mean is $\dfrac{482}{100}=\boxed{4.82}$.
